\chapter{The cd-index and Stanley's Fact \#3}
\label{ch:cdindex}

This chapter formalizes Stanley EC1 \S1.6.3, pp.~57--60: the cd-letter
classification of internal vertices, the cd-monomial $\PhiW{w}$, the
characteristic monomial $\charmono{S}$ of a descent set, and the
multiset identity (\emph{Fact \#3}, p.~58)
\[
  \sum_{v \in \Mclass{w}} \charmono{\DescSet{v}}
  \;=\;
  \PhiW{w}(a + b,\; ab + ba).
\]
The substantive support theorem \emph{$\PhiW{w} = \sum_{\omegaSet{X}
\supseteq S_w} \charmono{X}$} (Stanley's identity inside the proof of
Prop 1.6.4, p.~60--61) is proved here as
\rocq{mathcomp_eulerian.psi_cdindex_support.phi_w_support_general}.

% ============================================================
\section{The cd-letter alphabet}
% ============================================================

\begin{definition}[cd-letter alphabet]
  \label{def:cde}
  \rocq{mathcomp_eulerian.psi_cdindex_defs.cde}
  \rocqok
  $\code{cde} \;::=\; \mathtt{C\_letter} \;\mid\; \mathtt{D\_letter}
  \;\mid\; \mathtt{E\_letter}$.  These correspond to Stanley's
  cd-letters $\mathtt{c}$, $\mathtt{d}$ and the deleted endpoint
  $\mathtt{e}$.
\end{definition}

\begin{definition}[Vertex classification]
  \label{def:classify_vertex_cde}
  \uses{def:cde, def:is_internal, def:has_left_child}
  \rocq{mathcomp_eulerian.psi_cdindex_defs.classify_vertex_cde}
  \rocqok
  Position $i$ in $w$ is classified as:
  \[
    \mathrm{classify\_vertex\_cde}(i, w)
    \;=\;
    \begin{cases}
      \mathtt{E\_letter} & \text{if $i$ is not internal,} \\
      \mathtt{D\_letter} & \text{if $i$ is internal with left child,} \\
      \mathtt{C\_letter} & \text{if $i$ is internal without left child.}
    \end{cases}
  \]
\end{definition}

\begin{definition}[Unfiltered and cd-monomials]
  \label{def:phi_w}
  \uses{def:classify_vertex_cde}
  \rocq{mathcomp_eulerian.psi_cdindex_defs.phi'_w}
  \rocq{mathcomp_eulerian.psi_cdindex_defs.phi_w}
  \rocqok
  $\PhiprimeW{w} = [\,\mathrm{classify\_vertex\_cde}(i, w) : i \in \{0,
  \dots, |w|-1\}\,]$ is the unfiltered list of cd-letters by position.
  $\PhiW{w}$ is $\PhiprimeW{w}$ with $\mathtt{E\_letter}$s filtered out
  — Stanley's $\Phi_w$, the cd-monomial associated with $w$ (EC1
  pp.~58--59).
\end{definition}

\begin{theorem}[Stanley Fact \#1, well-definedness]
  \label{thm:phi_w_apply_psis}
  \uses{def:phi_w, def:apply_psis, lem:shape_apply_psis}
  \rocq{mathcomp_eulerian.psi_cdindex_core.phi_w_apply_psis}
  \rocqok
  For every $\mathrm{ss}$ and every uniq $w$,
  $\PhiW{\code{apply\_psis}(\mathrm{ss})(w)} = \PhiW{w}$.

  Equivalently: $\PhiW{w}$ depends only on the M-equivalence class
  $\Mclass{w}$ — equivalently, only on $\Mtree{w}$ as an unlabelled
  tree.  This is the well-definedness clause of Stanley's
  \emph{Fact \#1} (EC1 p.~57).  Internally it follows from the
  shape-invariance lemmas in
  \cref{lem:shape_apply_psis}.
\end{theorem}

% ============================================================
\section{The cd $\to$ ab expansion}
% ============================================================

\begin{definition}[Bit-string expansion of cd-monomials]
  \label{def:expand_cde}
  \uses{def:cde}
  \rocq{mathcomp_eulerian.psi_cdindex_defs.expand_cde}
  \rocqok
  $\code{expand\_cde}\,m \in \mathrm{seq}\,(\mathrm{seq}\,\mathrm{bool})$
  substitutes $\mathtt{c} \mapsto a + b$, $\mathtt{d} \mapsto ab + ba$,
  $\mathtt{e} \mapsto 1$ in the cd-monomial $m$:
  \begin{align*}
    \code{expand\_cde}\,[\,] &= [\,[\,]\,]\\
    \code{expand\_cde}\,(\mathtt{C} :: \mathit{rest}) &=
      [\,\mathtt{false} :: t : t \in \code{expand\_cde}\,\mathit{rest}\,]\\
    &\quad +\!\!+\;
      [\,\mathtt{true} :: t : t \in \code{expand\_cde}\,\mathit{rest}\,]\\
    \code{expand\_cde}\,(\mathtt{D} :: \mathit{rest}) &=
      [\,\mathtt{false}::\mathtt{true}::t : t \in \code{expand\_cde}\,\mathit{rest}\,]\\
    &\quad +\!\!+\;
      [\,\mathtt{true}::\mathtt{false}::t : t \in \code{expand\_cde}\,\mathit{rest}\,]\\
    \code{expand\_cde}\,(\mathtt{E} :: \mathit{rest}) &=
      \code{expand\_cde}\,\mathit{rest}.
  \end{align*}
  Each element of the result is a bit-string representing a descent
  pattern; the total length is $2^{c} \cdot 2^{d}$ where $c, d$ are
  the C- and D-counts of $m$.
\end{definition}

\begin{definition}[Sequence-level descent test]
  \label{def:is_descent_seq}
  \rocq{mathcomp_eulerian.psi_descent_v2.is_descent_seq}
  \rocqok
  For a sequence $w$ and $k \in \NN$,
  $\code{is\_descent\_seq}(w, k) = (w_{k+1} < w_k)$ — the
  sequence-level analogue of the permutation-level
  \cref{def:descent_set}.
\end{definition}

\begin{definition}[Characteristic monomial]
  \label{def:char_mono}
  \uses{def:is_descent_seq}
  \rocq{mathcomp_eulerian.psi_cdindex_defs.char_mono}
  \rocqok
  $\charmono{w} \;=\; [\,\code{is\_descent\_seq}(w, k) : k \in \{0,
  \dots, |w| - 2\}\,]$ is the bit-vector of descent positions.

  For a permutation $\sigma$, this matches Stanley's $\charmono{S} =
  e_1 e_2 \cdots e_{n-1}$ (eq.~1.60, p.~58) with $e_i = a$ if
  $i \notin S$ and $e_i = b$ if $i \in S$.  Our representation uses
  $\mathtt{false} \mapsto a$ and $\mathtt{true} \mapsto b$.
\end{definition}

% ============================================================
\section{Stanley's identity for $\PhiW{w}$}
% ============================================================

The next theorem is the load-bearing intermediate step in Stanley's
proof of Proposition 1.6.4 (p.~60--61).  It says that the bit-strings
in $\code{expand\_cde}\,(\PhiW{w})$ are exactly those whose induced
descent pattern is compatible with $S_w$ (the position set of D-letters
in $\PhiW{w}$, see \cref{ch:omega} and
Chapter~\ref{ch:propositions}).

\begin{theorem}[Stanley's $\PhiW{w} = \sum_{\omegaSet{X} \supseteq S_w} \charmono{X}$]
  \label{thm:phi_w_support_general}
  \uses{def:phi_w, def:expand_cde, def:char_mono, def:apply_psis}
  \rocq{mathcomp_eulerian.psi_cdindex_support.phi_w_support_general}
  \rocqok
  For every uniq $w$ with $|w| \geq 2$ and every bit-string $X$ with
  $|X| = |w| - 1$:
  \[
    X \in \code{expand\_cde}\,(\PhiW{w})
    \;\Longleftrightarrow\;
    \forall k \in S_w,\;\;
       k \in \omegaSet{\,\text{descent positions of }X\,}.
  \]
  This formalizes Stanley's identity at the top of EC1 p.~61
  inside the proof of Prop 1.6.4.

  The supporting bit-level definitions
  ($\code{cde\_width}$, $\code{cde\_total\_width}$,
  $\code{cde\_offset}$, $\code{D\_offsets}$,
  $\code{expand\_cde\_mem\_iff}$) live in
  \rocq{mathcomp_eulerian.psi_cdindex_support_defs}.
\end{theorem}

% ============================================================
\section{Stanley Fact \#3}
% ============================================================

\begin{theorem}[Stanley Fact \#3 (EC1 p.~58)]
  \label{thm:fact3}
  \uses{def:phi_w, def:expand_cde, def:char_mono, def:apply_psis, thm:psi_comm, thm:psi_involutive}
  \rocq{mathcomp_eulerian.psi_cdindex_support.fact3}
  \rocqok
  Stanley's statement (verbatim, p.~58):
  \begin{quote}
    \emph{Let $w \in \Sn$, and let $\Mclass{w}$ be the M-equivalence
    class containing $w$.  Then
    $\sum_{v \in \Mclass{w}} \charmono{\DescSet{v}} =
    \PhiW{w}(a + b,\; ab + ba)$.}
  \end{quote}
  Our formalization, with multisets represented as sorted seqs:
  \[
    \mathrm{sort}\!\left(\,\{\,\charmono{\code{apply\_psis}(\mathrm{ss})(w)} :
      \mathrm{ss} \subseteq \code{internal\_vertices}(w)\,\}\,\right)
    \;=\;
    \mathrm{sort}\!\left(\,\code{expand\_cde}\,(\PhiW{w})\,\right).
  \]

  Stanley uses Fact \#3 as a stepping stone to \emph{Theorem 1.6.3}
  (the ab-index of $\Sn$ factors through the cd-substitution); our
  $\code{fact3}$ matches this directly.

  \begin{proof}[Proof outline]
    The right-hand side is enumerated by Stanley's bit-string
    expansion (\cref{def:expand_cde}), and the left-hand side is
    enumerated by ranging over subsets of $\code{internal\_vertices}(w)$.
    Both have size $2^{\iota(w)}$ where $\iota(w)$ is the number of
    internal vertices.

    The map $\mathrm{ss} \mapsto \charmono{\code{apply\_psis}(\mathrm{ss})(w)}$
    is shown to be injective via a bit-level recovery argument (in
    \rocq{mathcomp_eulerian.psi_cdindex_core}): for each internal
    vertex $v$, one can read off $(v \in \mathrm{ss})$ from the bits
    of $\charmono{\code{apply\_psis}(\mathrm{ss})(w)}$ at positions
    owned by $v$.  Disjointness of the bit-positions across
    distinct internal vertices follows from
    \rocq{mathcomp_eulerian.psi_cdindex_core.LR_pred_is_endpoint}
    (D-vertex predecessors are non-internal).  Equality of the two
    multisets then follows from injectivity plus equal cardinality.
  \end{proof}
\end{theorem}
